% GENERATED FILE. Edit canonical lessons, then run npm run generate:books.
% canonical-source-hash: fnv1a64:669301f5513f8049
% canonical-lessons: HI-C109-hansi, HI-C109-ansu, HI-C109-dukh, HI-C109-akela, HI-C109-thanda

\chapter{Laughter and Tears}
\label{ch:hi-laughter-and-tears}
\hlchaptermodality{109}

\begin{chapteropening}
\textbf{By the end of this chapter:} \emph{I can talk about laughter and tears, sorrow, being alone, and something cold to drink.}

Complete HI-C109-thanda: moods, and one plain comfort.
\end{chapteropening}

\section[\texorpdfstring{\hi{हँसी} (hãsī)}{हँसी (hãsī)}]{\hi{हँसी} (hãsī) — laughter}

\label{lesson:HI-C109-hansi}



\begin{quote}
Before the new one: say the Hindi for troubled, then the Hindi for worry.
\end{quote}

\subsection*{You'll want to know: \hi{हँसी}}
\textbf{\hi{हँसी}} (\emph{hãsī}) — \textquotedblleft{}laughter\textquotedblright{}. It is feminine.

\textbf{\hi{मुझे} \hi{हँसी} \hi{आ} \hi{गई}} — \textquotedblleft{}I burst out laughing\textquotedblright{}, literally \textquotedblleft{}laughter came to me\textquotedblright{}, like anger.

\begin{cousinweb}
\textbf{\hi{हँसी}} (\emph{hãsī}) is a tadbhava, a word inherited from Sanskrit and reshaped by sound change. It comes from \emph{hasita}, \textquotedblleft{}laughter\textquotedblright{} (CDIAL 14023).
\end{cousinweb}

\begin{grammarlens}[title={What you've built}]
Laughter, which comes to you like anger does.
\end{grammarlens}

\subsection*{Your turn}
\begin{itemize}
\raggedright
  \item \emph{Say it:} \emph{hãsī}
  \item \emph{Say it:} \emph{hãsī}, once more
  \item \emph{Say it:} say \emph{mujhe hãsī ā gaī}
  \item \emph{Recall:} say the Hindi for troubled, then the Hindi for worry, then say \emph{hãsī} again
\end{itemize}

\begin{tcolorbox}[breakable,colback=teal!4,colframe=teal!35!black,title={Before you move on}]
What does \hi{हँसी} mean? (\textquotedblleft{}Laughter\textquotedblright{}.) Say it once more.
\end{tcolorbox}
\section[\texorpdfstring{\hi{आँसू} (ā̃sū)}{आँसू (ā̃sū)}]{\hi{आँसू} (ā̃sū) — tears}

\label{lesson:HI-C109-ansu}



\begin{quote}
Before the new one: say the Hindi for worry, then the Hindi for laughter.
\end{quote}

\subsection*{You'll want to know: \hi{आँसू}}
\textbf{\hi{आँसू}} (\emph{ā̃sū}) — \textquotedblleft{}tears\textquotedblright{}, a teardrop.

\textbf{\hi{आँख}} is the eye you know; \textbf{\hi{आँसू}} is what it sheds. Both open on the same nasal \textbf{\hi{आँ}}.

\begin{cousinweb}
\textbf{\hi{आँसू}} (\emph{āṁsū}) is a tadbhava, a word inherited from Sanskrit and reshaped by sound change. It comes from \emph{aśru}, \textquotedblleft{}tear\textquotedblright{} (CDIAL 919). CDIAL notes that the word took a nasal vowel early, as the word for \textquotedblleft{}eye\textquotedblright{} did.
\end{cousinweb}

\begin{grammarlens}[title={What you've built}]
Tears, beside the eye.
\end{grammarlens}

\subsection*{Your turn}
\begin{itemize}
\raggedright
  \item \emph{Say it:} \emph{ā̃sū}
  \item \emph{Say it:} \emph{ā̃sū}, once more
  \item \emph{Say it:} say \emph{ā̃sū}, then \emph{ā̃kh}
  \item \emph{Recall:} say the Hindi for worry, then the Hindi for laughter, then say \emph{ā̃sū} again
\end{itemize}

\begin{tcolorbox}[breakable,colback=teal!4,colframe=teal!35!black,title={Before you move on}]
What does \hi{आँसू} mean? (\textquotedblleft{}Tears\textquotedblright{}.) Say it once more.
\end{tcolorbox}
\section[\texorpdfstring{\hi{दुख} (dukh)}{दुख (dukh)}]{\hi{दुख} (dukh) — sorrow, grief}

\label{lesson:HI-C109-dukh}



\begin{quote}
Before the new one: say the Hindi for laughter, then the Hindi for tears.
\end{quote}

\subsection*{You'll want to know: \hi{दुख}}
\textbf{\hi{दुख}} (\emph{dukh}) — \textquotedblleft{}sorrow, grief\textquotedblright{}.

\textbf{\hi{मुझे} \hi{दुख} \hi{है}} — \textquotedblleft{}I am sorry\textquotedblright{}, said in sympathy, when someone tells you bad news. It is not the \textbf{\hi{माफ़} \hi{कीजिए}} of an apology.

\begin{cousinweb}
\textbf{\hi{दुख}} (\emph{dukh}) is a tadbhava, a word inherited from Sanskrit and reshaped by sound change. It comes from \emph{duḥkha}, \textquotedblleft{}difficulty, pain\textquotedblright{} (CDIAL 6375). CDIAL notes that \emph{duḥkha} was formed after \emph{sukha}, \textquotedblleft{}ease\textquotedblright{}.
\end{cousinweb}

\begin{grammarlens}[title={What you've built}]
The sympathy sorry, apart from the apology sorry.
\end{grammarlens}

\subsection*{Your turn}
\begin{itemize}
\raggedright
  \item \emph{Say it:} \emph{dukh}
  \item \emph{Say it:} \emph{dukh}, once more
  \item \emph{Say it:} say \emph{mujhe dukh hai}
  \item \emph{Recall:} say the Hindi for laughter, then the Hindi for tears, then say \emph{dukh} again
\end{itemize}

\begin{tcolorbox}[breakable,colback=teal!4,colframe=teal!35!black,title={Before you move on}]
What does \hi{दुख} mean? (\textquotedblleft{}Sorrow, grief\textquotedblright{}.) Say it once more.
\end{tcolorbox}
\section[\texorpdfstring{\hi{अकेला} (akelā)}{अकेला (akelā)}]{\hi{अकेला} (akelā) — alone}

\label{lesson:HI-C109-akela}



\begin{quote}
Before the new one: say the Hindi for tears, then the Hindi for sorrow.
\end{quote}

\subsection*{You'll want to know: \hi{अकेला}}
\textbf{\hi{अकेला}} (\emph{akelā}) — \textquotedblleft{}alone\textquotedblright{}. A woman says \textbf{\hi{अकेली}}.

\textbf{\hi{मैं} \hi{अकेला} \hi{हूँ}} — \textquotedblleft{}I am alone\textquotedblright{}. It opens on \textbf{\hi{एक}}, \textquotedblleft{}one\textquotedblright{}, worn down.

\begin{cousinweb}
\textbf{\hi{अकेला}} (\emph{akelā}) is a tadbhava, a word inherited and reshaped by sound change. It comes from \emph{*ekkala}, \textquotedblleft{}alone\textquotedblright{}, a form CDIAL reconstructs from its descendants (CDIAL 2506). It is built on \emph{eka}, \textquotedblleft{}one\textquotedblright{}.
\end{cousinweb}

\begin{grammarlens}[title={What you've built}]
Alone, from the word for one.
\end{grammarlens}

\subsection*{Your turn}
\begin{itemize}
\raggedright
  \item \emph{Say it:} \emph{akelā}
  \item \emph{Say it:} \emph{akelā}, once more
  \item \emph{Say it:} say \emph{maiṁ akelā hūṁ}, then \emph{maiṁ akelī hūṁ}
  \item \emph{Recall:} say the Hindi for tears, then the Hindi for sorrow, then say \emph{akelā} again
\end{itemize}

\begin{tcolorbox}[breakable,colback=teal!4,colframe=teal!35!black,title={Before you move on}]
What does \hi{अकेला} mean? (\textquotedblleft{}Alone\textquotedblright{}.) Say it once more.
\end{tcolorbox}
\section[\texorpdfstring{\hi{ठंडा} (ṭhaṇḍā)}{ठंडा (ṭhaṇḍā)}]{\hi{ठंडा} (ṭhaṇḍā) — cold, cool}

\label{lesson:HI-C109-thanda}



\begin{quote}
Before the new one: say the Hindi for sorrow, then the Hindi for alone.
\end{quote}

\subsection*{You'll want to know: \hi{ठंडा}}
\textbf{\hi{ठंडा}} (\emph{ṭhaṇḍā}) — \textquotedblleft{}cold, cool\textquotedblright{}, of a thing.

\textbf{\hi{ठंडा} \hi{पानी}} — \textquotedblleft{}cold water\textquotedblright{}, the first comfort offered to a guest on a hot day.

\begin{cousinweb}
\textbf{\hi{ठंडा}} (\emph{ṭhaṇḍā}) comes from a reconstructed \emph{*ṭhaṇḍha}, \textquotedblleft{}cold\textquotedblright{}, which CDIAL lists under Sanskrit \emph{stabdha}, \textquotedblleft{}stiff, solidified\textquotedblright{}, tracing the sense from \textquotedblleft{}firm\textquotedblright{} through \textquotedblleft{}sluggish\textquotedblright{} to \textquotedblleft{}cold\textquotedblright{} (CDIAL 13676). DEDR compares Dravidian words for \textquotedblleft{}cool\textquotedblright{}, such as Tamil \emph{taṇ} (DEDR 3045).
\end{cousinweb}

\begin{grammarlens}[title={What you've built}]
Laughter, tears, sorrow, alone, and something cold to drink.
\end{grammarlens}

\subsection*{Your turn}
\begin{itemize}
\raggedright
  \item \emph{Say it:} \emph{ṭhaṇḍā}
  \item \emph{Say it:} \emph{ṭhaṇḍā}, once more
  \item \emph{Say it:} say \emph{ṭhaṇḍā pānī}
  \item \emph{Recall:} say the Hindi for sorrow, then the Hindi for alone, then say \emph{ṭhaṇḍā} again
\end{itemize}

\begin{tcolorbox}[breakable,colback=teal!4,colframe=teal!35!black,title={Before you move on}]
What does \hi{ठंडा} mean? (\textquotedblleft{}Cold, cool\textquotedblright{}.) Say it once more.
\end{tcolorbox}
